\section*{Spec 1: A Formal Model for Deterministic Filesystem Space Analysis}

\subsection*{Introduction}

We present a formal model for deterministic disk space analysis over a finite filesystem.
The objective is to define the mathematical structures, invariants, and aggregation
properties that a correct implementation must satisfy, independent of operating system or implementation language.

\subsection*{Filesystem Model}

\begin{definition}[Filesystem Graph]
A filesystem is a finite directed multigraph
\[
\mathcal{F} = (V, E, \tau, \mu),
\]
where $V$ is a finite set of nodes, $E \subseteq V \times V$ is a set of directed edges,
$\tau : V \to T$ assigns a type to each node, and
$\mu : V \to \mathbb{N}$ assigns a non-negative size measure.
The type set is $T = \{\mathrm{file}, \mathrm{dir}, \mathrm{symlink}, \mathrm{special}\}$.
\end{definition}

We distinguish a subset $E_d \subseteq E$ representing directory containment.
The induced subgraph $(V, E_d)$ is acyclic when restricted to nodes of type $\mathrm{dir}$ and $\mathrm{file}$.

\begin{definition}[Containment Forest]
Let $V_0 = \{ v \in V \mid \tau(v) \in \{\mathrm{file}, \mathrm{dir}\} \}$.
Then $(V_0, E_d)$ forms a finite forest.
\end{definition}

\subsection*{Aggregation Semantics}

Let $r \in V$ be a designated root.

\begin{definition}[Recursive Aggregation]
Define $A : V \to \mathbb{N}$ recursively by
\[
A(v) =
\begin{cases}
\mu(v), & \tau(v) = \mathrm{file}, \\
\sum_{(v,u) \in E_d} A(u), & \tau(v) = \mathrm{dir}.
\end{cases}
\]
\end{definition}

\begin{proposition}
For any directory $v$,
\[
A(v) = \sum_{u \in \mathrm{Desc}_f(v)} \mu(u).
\]
\end{proposition}

\begin{proof}
Proceed by induction on directory height.

If $v$ contains only files, then
\[
A(v) = \sum_{(v,u)\in E_d} \mu(u),
\]
which equals the sum over $\mathrm{Desc}_f(v)$.

Suppose the statement holds for all strict subdirectories of $v$.
Then
\[
A(v) = \sum_{(v,u)\in E_d} A(u).
\]

Each $u$ is either a file or a directory.
If $u$ is a file, $A(u)=\mu(u)$.
If $u$ is a directory, by the inductive hypothesis,
\[
A(u) = \sum_{w \in \mathrm{Desc}_f(u)} \mu(w).
\]

Therefore
\[
A(v) = \sum_{u \in \mathrm{Desc}_f(v)} \mu(u).
\]
\end{proof}

\subsection*{Hard Link Equivalence}

Let $\sim$ be an equivalence relation on file nodes representing shared inodes.

\begin{definition}[Inode Quotient]
Let $\tilde{V} = V / \sim$ and define
\[
\tilde{\mu}([v]) = \mu(v).
\]
Aggregation must be computed over equivalence classes rather than individual representatives.
\end{definition}

\begin{theorem}[No Double Counting]
If aggregation is performed over $\tilde{V}$, then each physical storage unit contributes exactly once to $A(r)$.
\end{theorem}

\begin{proof}
Each equivalence class corresponds to a unique physical allocation.
Since $\tilde{\mu}$ is defined on classes and summation is over distinct classes,
duplication is eliminated.
\end{proof}

\subsection*{Traversal Correctness}

\begin{definition}[Traversal Function]
A traversal is a function
\[
T : V \to \{0,1\}
\]
indicating whether a node has been visited.
\end{definition}

\begin{proposition}[Termination]
If traversal maintains a visited set and $V$ is finite, then traversal terminates even in the presence of symbolic link cycles.
\end{proposition}

\begin{proof}
Each node enters the visited set at most once.
Since $|V|<\infty$, recursion depth is finite.
\end{proof}

\subsection*{Metric Algebra}

Let $(\mathbb{N}, +, 0)$ denote the additive monoid of non-negative integers.

\begin{definition}[Monoidal Fold]
Directory aggregation is a fold over $(V_0,E_d)$ with respect to the monoid $(\mathbb{N}, +, 0)$.
\end{definition}

\begin{proposition}
Aggregation is associative and independent of traversal order.
\end{proposition}

\begin{proof}
Associativity and commutativity of $+$ imply order independence.
\end{proof}

\subsection*{Determinism}

\begin{theorem}[Deterministic Output]
If two filesystem graphs $\mathcal{F}_1$ and $\mathcal{F}_2$ are isomorphic
and preserve $\mu$ and $\tau$, then aggregation outputs are equal.
\end{theorem}

\begin{proof}
Aggregation depends only on graph structure and measure $\mu$.
Graph isomorphism preserving these structures yields identical sums.
\end{proof}

\subsection*{Complexity Bound}

Let $n = |V|$.

\begin{proposition}
A complete traversal with aggregation requires $O(n)$ time.
\end{proposition}

\begin{proof}
Each node and edge is processed at most once.
\end{proof}

\subsection*{Memory Bound}

Let $h$ denote maximum directory depth.

\begin{proposition}
A depth-first traversal requires $O(h)$ auxiliary memory.
\end{proposition}

\begin{proof}
At any time, the call stack contains at most one node per depth level.
\end{proof}

\subsection*{Conclusion}

The Drive Space Analyzer is formally characterized as a deterministic monoidal aggregation
over a finite containment forest derived from a filesystem graph.
Correctness requires quotienting by inode equivalence, ensuring termination under cyclic symbolic references,
and preserving traversal-order invariance.